\section{Технико-экономическая оценка и анализ эксплуатационных рисков}

\subsection{Подход к технико-экономической оценке}

Экономическая часть для рассматриваемой платформы должна оценивать не только
стоимость написания кода, но и эксплуатационный эффект архитектурных решений.
Система уведомлений является инфраструктурным компонентом, поэтому её ценность
проявляется прежде всего в снижении операционных потерь: потерь уведомлений,
затрат на ручную диагностику, стоимости инцидентов и времени восстановления.

В рамках данной ВКР корректно использовать качественную технико-экономическую
оценку. Для строгого количественного расчёта по нормативной вузовской методике
необходимы утверждённые исходные показатели: ставки исполнителей, нормативы
трудоёмкости, тарифы на инфраструктуру, стоимость владения внешними шлюзами,
регламентированные коэффициенты накладных расходов и принятая кафедрой схема
приведения затрат.

\assumption{В текущем репозитории и сопроводительных материалах отсутствуют
утверждённые кафедрой числовые исходные данные для полного калькуляционного
расчёта себестоимости и экономического эффекта. Поэтому раздел выполнен как
качественная технико-экономическая оценка с явным анализом факторов стоимости,
эффекта и рисков, без вымышленных денежных величин.}

\subsection{Основные статьи затрат}

Для платформы мультиканальной доставки уведомлений можно выделить четыре группы
затрат.

\subsubsection{Затраты на разработку}

В эту группу входят:
\begin{enumerate}
    \item проектирование доменной модели события, запуска доставки и канальной доставки;
    \item реализация gRPC-контрактов и Kafka-пайплайна;
    \item создание миграций PostgreSQL и Redis-модели профиля;
    \item разработка механизмов идемпотентности, outbox/inbox и повторных попыток;
    \item реализация sender-сервисов и интеграции с историей.
\end{enumerate}

На начальном этапе такие затраты выше, чем у упрощённого монолитного сервиса с
прямой синхронной отправкой. Однако это увеличение затрат компенсируется снижением
технического долга и уменьшением стоимости последующих изменений.

\subsubsection{Инфраструктурные затраты}

К инфраструктурной части относятся:
\begin{enumerate}
    \item вычислительные ресурсы прикладных сервисов;
    \item PostgreSQL как устойчивое хранилище доменного состояния;
    \item Redis как оперативное хранилище профилей;
    \item Kafka как межсервисный транспорт;
    \item Prometheus и Grafana как средства наблюдаемости.
\end{enumerate}

По сравнению с централизованной схемой количество компонентов возрастает, однако
каждый из них закрывает отдельный класс инженерных рисков: Kafka уменьшает
связанность и позволяет управлять асинхронной нагрузкой, Redis ускоряет доступ к
профилю получателя, а Prometheus и Grafana сокращают стоимость диагностики.

\subsubsection{Эксплуатационные затраты}

В эксплуатационные затраты входят:
\begin{enumerate}
    \item сопровождение контейнерного контура;
    \item контроль миграций схем данных;
    \item поддержка релизного процесса и стендов;
    \item мониторинг инцидентов и разбор проблемных доставок;
    \item администрирование секретов и политик безопасности.
\end{enumerate}

Именно на этой стадии архитектурные решения начинают давать экономический эффект.
Наличие истории, явных статусов и метрик сокращает затраты на поиск причины
проблемы и уменьшает число ручных операций при восстановлении обработки.

\subsection{Экономический эффект от архитектурных решений}

\subsubsection{Эффект от разделения жизненного цикла}

Разделение на событие, запуск доставки и отдельную канальную доставку позволяет
избежать целого класса ошибок, когда повторный запуск или частичная неуспешность
ломают общее состояние уведомления. С экономической точки зрения это снижает
стоимость сопровождения, поскольку каждая стадия имеет собственные таблицы,
собственные статусы и отдельную точку анализа.

\subsubsection{Эффект от outbox/inbox-подхода}

Паттерны outbox и inbox увеличивают сложность реализации, но уменьшают риск
тихой потери сообщений и трудоёмкость восстановления после сбоя. Для сервиса
уведомлений это особенно важно: стоимость одной потерянной транзакционной
коммуникации может значительно превышать стоимость дополнительной технической
логики в коде.

\subsubsection{Эффект от централизованного профиля получателя}

Выделение profile-consent как отдельного сервиса устраняет дублирование логики
проверки согласий, блокировок и каналов в sender-сервисах. Экономически это
снижает стоимость развития новых каналов и уменьшает риск рассинхронизации
бизнес-правил между разными отправителями.

\subsubsection{Эффект от наблюдаемости}

Метрики Kafka-публикации, listener-обработки и gRPC-вызовов создают прямой
эксплуатационный эффект: инженер может определить, где именно возникла проблема,
не прибегая к догадкам и трудоёмкому анализу всех журналов вручную. Это уменьшает
время локализации инцидента и тем самым снижает операционные затраты.

\subsection{Сравнительная оценка вариантов}

В таблице~\ref{tab:economic-comparison} приведено качественное сравнение двух
вариантов: упрощённого централизованного контура и реализованной микросервисной
архитектуры.

\begin{table}[h]
    \centering
    \footnotesize
    \renewcommand{\arraystretch}{1.15}
    \caption{Качественное сравнение архитектурных вариантов}
    \label{tab:economic-comparison}
    \begin{tabular}{|L{0.34\textwidth}|L{0.25\textwidth}|L{0.25\textwidth}|}
        \hline
        \textbf{Критерий} & \textbf{Централизованный сервис} & \textbf{Реализованная архитектура} \\
        \hline
        Начальная стоимость разработки & Ниже & Выше \\
        \hline
        Стоимость локализации сбоя & Выше & Ниже \\
        \hline
        Стоимость расширения новым каналом & Выше из-за сильной связанности & Ниже за счёт выделенных сервисных границ \\
        \hline
        Риск потери сообщения при сбое & Выше & Ниже благодаря outbox/inbox \\
        \hline
        Стоимость поддержки истории и аудита & Выше & Ниже за счёт отдельного history-сервиса \\
        \hline
        Предсказуемость поведения при повторе & Ниже & Выше благодаря идемпотентности и уникальным ограничениям \\
        \hline
    \end{tabular}
\end{table}

Таким образом, реализованное решение экономически оправдано не в логике
«минимальная цена входа», а в логике уменьшения будущих эксплуатационных и
организационных потерь.

\subsection{Анализ рисков информационной безопасности}

Для платформы уведомлений критичны не только экономические, но и
информационно-безопасностные риски, поскольку система оперирует данными о
получателях, каналами доставки и служебными идентификаторами.
Базовые меры защиты межсервисного взаимодействия в таких системах опираются на
шифрование транспортного канала и контролируемую сервисную аутентификацию
\cite{rfc8446,rfc6749,rfc7519}.

Ключевые риски и меры снижения представлены в
таблице~\ref{tab:security-risks}.

\begin{table}[h]
    \centering
    \footnotesize
    \renewcommand{\arraystretch}{1.15}
    \caption{Ключевые риски информационной безопасности}
    \label{tab:security-risks}
    \begin{tabular}{|L{0.28\textwidth}|L{0.32\textwidth}|L{0.26\textwidth}|}
        \hline
        \textbf{Риск} & \textbf{Последствие} & \textbf{Меры снижения} \\
        \hline
        Несанкционированный вызов фасада & Запуск ложных уведомлений, нарушение бизнес-процессов & Сервисная аутентификация, проверка client\_id, аудит вызовов \\
        \hline
        Перехват межсервисного трафика & Компрометация payload и служебных идентификаторов & TLS, сервисные сертификаты, сегментация сети \\
        \hline
        Утечка секретов шлюзов & Несанкционированная отправка через внешних провайдеров & Внешнее хранилище секретов, ротация, ограничение доступа \\
        \hline
        Чтение данных профиля из Redis без контроля доступа & Утечка контактов и признаков согласия & Закрытие прямого доступа, использование только через profile-consent \\
        \hline
        Избыточное логирование payload & Попадание персональных данных в журналы & Маскирование чувствительных полей, политика логирования, контроль retention \\
        \hline
    \end{tabular}
\end{table}

Отдельно необходимо учитывать, что эксплуатационная стоимость инцидента
безопасности для подобной платформы обычно выше стоимости превентивных мер. Даже
если меры защиты усложняют инфраструктуру, они экономически оправданы за счёт
снижения риска компрометации каналов и получательских данных.

\subsection{Влияние архитектуры на стоимость сопровождения}

С точки зрения сопровождения наиболее экономически значимыми свойствами системы
являются:
\begin{enumerate}
    \item отдельный сервис истории, позволяющий разбирать инциденты по фактам, а не по косвенным признакам;
    \item явные доменные статусы и таблицы попыток отправки;
    \item централизованный profile-consent, уменьшающий дублирование логики в sender-сервисах;
    \item возможность тестировать и дорабатывать отдельный контур без переписывания всей системы;
    \item наблюдаемость каждого этапа через метрики Kafka и gRPC.
\end{enumerate}

Именно эти свойства уменьшают полную стоимость владения в среднесрочной
перспективе. Архитектурная дисциплина требует больших усилий на старте, но затем
снижает цену сопровождения, потому что уменьшает хаотичность поведения системы.

\subsection{Выводы по разделу}

Выполненная технико-экономическая оценка показывает, что реализованная платформа
имеет более высокую стартовую сложность по сравнению с простым централизованным
сервисом, но экономически выигрывает на этапе эксплуатации. Основной эффект
достигается за счёт уменьшения потерь сообщений, снижения трудоёмкости
диагностики, упрощения расширения новыми каналами и повышения управляемости
межсервисного взаимодействия. Дополнительный вклад в экономическую целесообразность
вносит архитектура безопасности: затраты на защитные меры оправданы тем, что
предотвращают существенно более дорогие инциденты, связанные с утечкой данных и
несанкционированной доставкой.


